\section*{Bloque A. Fundamentos teóricos y conceptuales}

\noindent\textit{Preguntas 1 a 30. Evalúan la definición de monopolio, las barreras a la entrada, la maximización del beneficio, el ingreso marginal, la discriminación de precios, la pérdida de eficiencia y el monopolio natural.}

\q{Un monopolio se define como:}
\begin{opts}
    \item una empresa única sin sustitutos cercanos y con barreras a la entrada.
    \item muchas empresas atomizadas que ofrecen un producto totalmente homogéneo.
    \item dos competidores estratégicos que fijan precios de manera simultánea.
    \item una empresa precio-aceptante que enfrenta demanda horizontal infinita.
\end{opts}

\q{El monopolista es precio-decisor porque:}
\begin{opts}
    \item enfrenta una curva de demanda de mercado perfectamente elástica.
    \item enfrenta demanda de pendiente negativa y elige punto precio-cantidad.
    \item opera sin costos fijos ni costos marginales crecientes en planta.
    \item recibe una tarifa máxima regulada fijada directamente por el Estado.
\end{opts}

\q{Una barrera a la entrada es:}
\begin{opts}
    \item la libertad irrestricta de entrada y salida de nuevos competidores.
    \item control de insumo esencial, una patente o economías de escala.
    \item la presencia de un bien perfectamente homogéneo sin diferenciación.
    \item la disponibilidad de información completa para todos los consumidores.
\end{opts}

\q{La diferencia central entre monopolio y competencia perfecta es que:}
\begin{opts}
    \item el monopolio opera invariablemente con curvas de costos medios más altos.
    \item la competencia perfecta concentra toda la oferta en un único vendedor.
    \item el monopolio fija precio sobre la demanda; la firma competitiva lo toma dado.
    \item el monopolista no busca maximizar beneficios económicos sino cobertura.
\end{opts}

\q{La curva de demanda que enfrenta el monopolio es:}
\begin{opts}
    \item perfectamente horizontal al nivel competitivo.
    \item perfectamente vertical e insensible al precio.
    \item infinitamente elástica para cualquier cantidad.
    \item la demanda de mercado, con pendiente negativa.
\end{opts}

\q{Para el monopolista, el ingreso marginal es menor que el precio porque:}
\begin{opts}
    \item vende la producción marginal a precio regulado nulo.
    \item para vender más debe bajar el precio de toda unidad.
    \item carece de costos marginales positivos de producción.
    \item enfrenta una curva de demanda de mercado horizontal.
\end{opts}

\q{La condición de maximización del beneficio del monopolista es:}
\begin{opts}
    \item $P=CMg$.
    \item $IMg=CMe$.
    \item $IMg=CMg$.
    \item $P=CMe$.
\end{opts}

\q{El precio de monopolio se encuentra:}
\begin{opts}
    \item exactamente en el mínimo del costo total medio.
    \item igual al costo marginal en el tramo decreciente.
    \item sobre la demanda, por encima del costo marginal.
    \item estrictamente por debajo del costo marginal.
\end{opts}

\q{Los beneficios económicos del monopolio pueden persistir en el largo plazo porque:}
\begin{opts}
    \item existen barreras a la entrada.
    \item hay libre entrada de empresas.
    \item el producto es homogéneo.
    \item la demanda es horizontal.
\end{opts}

\q{El monopolio no tiene una curva de oferta bien definida porque:}
\begin{opts}
    \item no busca maximizar beneficios económicos en el largo plazo.
    \item su costo marginal permanece siempre estrictamente constante.
    \item el precio depende de la demanda, sin relación única precio-cantidad.
    \item produce exactamente la misma cantidad fija en todo periodo.
\end{opts}

\q{El índice de Lerner mide:}
\begin{opts}
    \item el nivel mínimo del costo medio.
    \item la elasticidad precio de la demanda.
    \item el beneficio económico contable.
    \item el poder de mercado: $(P-CMg)/P$.
\end{opts}

\q{La regla de precios del monopolio con elasticidad constante es:}
\begin{opts}
    \item $P = CMg \cdot (1 - 1/|\varepsilon|)$ para $|\varepsilon|>1$
    \item $P = CMe / (1 + 1/|\varepsilon|)$ para $|\varepsilon|>0$
    \item $P = CMg \cdot (1 + 1/|\varepsilon|)$ para $|\varepsilon|<1$
    \item $P = CMg / (1 - 1/|\varepsilon|)$ para $|\varepsilon|>1$
\end{opts}

\q{El monopolista produce en el tramo de demanda:}
\begin{opts}
    \item elástico, con $|\varepsilon|>1$.
    \item inelástico, con $|\varepsilon|<1$.
    \item donde $\varepsilon=0$.
    \item donde $\varepsilon=\infty$.
\end{opts}

\q{La discriminación de precios consiste en:}
\begin{opts}
    \item cobrar precios distintos por razones ajenas a diferencias de costos.
    \item cobrar exactamente el mismo precio lineal uniforme a todo comprador.
    \item fijar el precio unitario exactamente igual al costo marginal de oferta.
    \item abastecer exclusivamente a un único comprador bajo contrato privado.
\end{opts}

\q{En la discriminación de precios de primer grado:}
\begin{opts}
    \item se cobra un único precio uniforme lineal.
    \item se cobra un precio distinto por segmento.
    \item se ofrecen paquetes de bloques por volumen.
    \item se cobra a cada cliente su máxima disposición.
\end{opts}

\q{La discriminación de precios de segundo grado:}
\begin{opts}
    \item cobra a cada consumidor exactamente su disposición máxima a pagar.
    \item usa autoselección mediante descuentos o bloques de consumo.
    \item separa grupos exógenos según su elasticidad precio de demanda.
    \item exige permitir el libre arbitraje entre todos los compradores.
\end{opts}

\q{La discriminación de precios de tercer grado:}
\begin{opts}
    \item aplica un precio lineal único a toda la demanda.
    \item utiliza menús autoseleccionados por volumen total.
    \item cobra precios distintos por grupo según elasticidad.
    \item extrae el cien por ciento del excedente consumidor.
\end{opts}

\q{Para que la discriminación de precios sea posible se requiere:}
\begin{opts}
    \item poder de mercado, segmentación y ausencia de arbitraje.
    \item estructura de competencia perfecta con libre entrada.
    \item libre entrada y salida sin costos hundidos relevantes.
    \item fijación del precio exactamente igual al costo marginal.
\end{opts}

\q{La discriminación de precios:}
\begin{opts}
    \item aumenta invariablemente el bienestar agregado total.
    \item puede subir o bajar el bienestar según cobertura y asignación.
    \item reduce necesariamente el bienestar bajo cualquier demanda.
    \item resulta completamente neutral respecto al bienestar social.
\end{opts}

\q{La pérdida de eficiencia del monopolio se mide por:}
\begin{opts}
    \item el monto total de beneficios extraordinarios del monopolio.
    \item el excedente neto retenido por los consumidores inframarginales.
    \item el triángulo entre demanda y costo marginal por unidades omitidas.
    \item la suma acumulada de los costos fijos hundidos no recuperables.
\end{opts}

\q{El monopolio genera pérdida de eficiencia porque:}
\begin{opts}
    \item produce una cantidad superior a la competitiva.
    \item fija $P>CMg$ y produce menos que el óptimo social.
    \item fija el precio exactamente igual al costo marginal.
    \item obtiene beneficios económicos nulos a largo plazo.
\end{opts}

\q{Un monopolio natural se caracteriza por:}
\begin{opts}
    \item costos medios crecientes en todo nivel productivo.
    \item una curva de oferta competitiva perfectamente horizontal.
    \item ausencia de barreras legales en competencia perfecta.
    \item costos medios decrecientes en el rango de demanda.
\end{opts}

\q{En un monopolio natural, fijar $P=CMg$ genera:}
\begin{opts}
    \item beneficios extraordinarios.
    \item pérdidas, porque $CMg<CMe$.
    \item beneficio económico cero.
    \item ausencia de pérdidas.
\end{opts}

\q{Una forma de regular un monopolio natural es:}
\begin{opts}
    \item fijar precio tope, tarifa en dos partes o subsidiar pérdidas.
    \item suprimir por completo las barreras legales permitiendo duplicar redes.
    \item establecer por mandato un precio uniforme igual a cero soles.
    \item prohibir toda nueva inversión de capital en ampliación de capacidad.
\end{opts}

\q{Una tarifa en dos partes consiste en:}
\begin{opts}
    \item un precio único proporcional por cada unidad.
    \item exclusivamente un cargo fijo mensual uniforme.
    \item un cargo fijo más un precio por unidad.
    \item un impuesto correctivo pigouviano específico.
\end{opts}

\q{El monopolio viola la eficiencia asignativa porque:}
\begin{opts}
    \item $P=CMg$ maximizando el bienestar agregado.
    \item produce una cantidad mayor a la eficiente.
    \item carece de costos marginales de producción.
    \item $P>CMg$, violando el Primer Teorema.
\end{opts}

\q{El monopolio puede ser productivamente ineficiente porque:}
\begin{opts}
    \item opera invariablemente en el mínimo del costo total medio.
    \item carece por completo de costos fijos de infraestructura.
    \item carece de incentivo a minimizar costos bajo barreras.
    \item enfrenta una curva de demanda perfectamente horizontal.
\end{opts}

\q{La ineficiencia X se refiere a:}
\begin{opts}
    \item costos excesivos por falta de presión competitiva.
    \item pérdida de excedente por un precio igual al costo marginal.
    \item discriminación perfecta.
    \item economía de escala.
\end{opts}

\q{Una política pública frente al monopolio puede ser:}
\begin{opts}
    \item suprimir por decreto la demanda residencial del bien.
    \item fijar precios regulados por encima del óptimo monopólico.
    \item prohibir el consumo industrial en el mercado relevante.
    \item regular, subastar concesiones, competencia o gestión pública.
\end{opts}

\q{Un precio máximo regulado en un monopolio natural, fijado en el costo medio, permite:}
\begin{opts}
    \item obtener beneficios económicos extraordinarios positivos.
    \item incurrir inevitablemente en pérdidas financieras crónicas.
    \item cubrir costos sin precio monopólico, sin primer mejor.
    \item suprimir totalmente la demanda del segmento residencial.
\end{opts}
